\section{Conclusion}

This paper's contribution divides cleanly into what is proved outright and what
remains conditional, and the two should not be conflated.

\textbf{Proved unconditionally.} The Algebraic Defect Framework itself ---
the root/fiber decomposition of the external boundary into a term linear in
$(n-k)$ minus a topology-dependent correction (Section~\ref{sec:defect}), and
the subadditivity lemma underlying it --- is established outright and
mechanized in Lean~4, independent of either open proposition below. On top of
it, the Fixed-Topology Asymptotic Penalty (Theorem~\ref{thm:penalty}) is an
unconditional theorem: for every fixed sub-optimal topology, the Hamming ball's
boundary advantage grows without bound as $n-k \to \infty$, which Cheng,
Lipt\'ak, and Tian's exhaustive small-case computations
\cite{cheng2022extraconnectivity} did not establish in general. The partial
defect-spectrum classification of Section~\ref{sec:fracture} (the Topological
Funnel, its primary and secondary fractures, and Theorem~\ref{thm:classification})
is likewise proved on paper from Harper's edge-isoperimetric theorem, not
merely observed computationally.

\textbf{Conditional.} Combining these unconditional results into a single
exact closed-form boundary formula valid within the hypercube embedding
regime ($R \le 2^m$) requires the Restricted Boundary Inequality
(Proposition~\ref{prop:restricted_lower_bound}), which holds for $m \le 4$ by
computational verification but remains open for general $m$.
A second, independent
gap --- the
boundary-to-extraconnectivity reduction relating minimum-boundary sets to
exact $(R{-}1)$-extraconnectivity --- is likewise open and is not addressed by
the Lean formalization, whose capstone theorem is confined to external-neighbor
sets of $R$-element subsets.

\textbf{Net assessment.} Relative to \cite{cheng2022extraconnectivity}'s exact
computation of small-$g$ extraconnectivity cases ($g \le 6$), this paper
establishes a conditional closed-form framework that subsumes those results
within the embeddable regime. For $m \le 4$ (covering all of the paper's
canonical examples), the framework is computationally verified. Closing the
Restricted Boundary Inequality for general $m$, and the
boundary-to-extraconnectivity reduction, remain the principal open problems.

\paragraph{Conjugate geometry and local induction.}
The computational diagnostics expose a persistent trade-off between the
root-incidence defect $D(V)$ and the cross-collision multiplicity
$\sum_{w:\,\mu_V(w)\ge2}\mu_V(w)$. Sparse Star configurations have small
defect and large collision capacity, while dense Hamming-ball configurations
have the opposite profile. These quantities therefore cannot be maximized
independently in a valid extremal argument. The exact coordinate-averaged
squeeze nevertheless fails on the full Star $S_4\subset A(8,4)$: its child
potential sum is $460$, its geometric penalty is $224$, and the arithmetic
budget is $4P(17)=676$, giving the exact deficit $684-676=8$. This records a
failure of a strictly local payment rule, not a proof of a replacement
theorem; the remaining boundary problem requires a coupled argument or a
different extremal description.
